\chapter{Fundamentação Teórica}
\label{chap:fundamentacao}

\section{Inteligência Artificial e Aprendizado de Máquina}
A Inteligência Artificial (IA) consolidou-se como um campo multidisciplinar da ciência da computação cujo escopo principal é a simulação de processos cognitivos humanos por meio de sistemas computacionais. No âmbito da IA, o Aprendizado de Máquina (\textit{Machine Learning}) destaca-se como o paradigma predominante, caracterizado pelo desenvolvimento de algoritmos que otimizam seu desempenho na execução de uma tarefa específica a partir da exposição empírica a dados históricos \cite{goodfellow2016deep}. 

Com o advento do Aprendizado Profundo (\textit{Deep Learning}), que se baseia em redes neurais artificiais com múltiplas camadas de neurônios ocultos, tornou-se viável processar representações de dados de alta dimensionalidade sem a necessidade de extração manual de características (\textit{handcrafted features}). Essas arquiteturas realizam a extração hierárquica de padrões de forma automática, sendo fundamentais para o avanço do processamento de imagens e sinais biológicos.

\section{Visão Computacional e Detecção de Objetos}
A visão computacional representa uma subárea da inteligência artificial cujo propósito é emular o sistema visual humano, permitindo que máquinas processem, analisem e compreendam imagens digitais e vídeos bidimensionais \cite{goodfellow2016deep}. Os avanços na área estão fortemente correlacionados à evolução das Redes Neurais Convolucionais (CNNs), cuja arquitetura utiliza filtros matemáticos locais compartilhados para capturar dependências espaciais e invariância de translação em imagens.

A detecção de objetos é uma tarefa clássica e de alta complexidade que unifica dois problemas fundamentais: a classificação de objetos (identificar ``o que'' está presente na imagem) e a localização espacial (determinar ``onde'' cada objeto está localizado). Essa localização é representada geometricamente por caixas delimitadoras (\textit{bounding boxes}), definidas por suas coordenadas espaciais e pelas respectivas probabilidades de classe \cite{redmon2016you}. A detecção multiclasse de objetos é a base tecnológica para sistemas de navegação autônoma, controle de tráfego, automação médica e monitoramento industrial de conformidades.

\section{Evolução e Arquitetura YOLO (\textit{You Only Look Once})}
Tradicionalmente, os sistemas de detecção de objetos baseavam-se em duas etapas (\textit{two-stage detectors}), como a família R-CNN (R-CNN, Fast R-CNN e Faster R-CNN). Esses sistemas realizam primeiro a geração de propostas de regiões de interesse na imagem e, em seguida, aplicam classificadores nessas regiões individuais. Embora apresentem alta precisão, tais modelos possuem elevado custo computacional, o que inviabiliza seu uso em aplicações de tempo real.

O paradigma YOLO (\textit{You Only Look Once}), proposto originalmente por \citeonline{redmon2016you}, revolucionou a área ao unificar a detecção em uma única etapa (\textit{one-stage detector}). O YOLO reformula o problema de detecção como uma regressão direta: a imagem inteira é processada por uma única rede convolucional, que prediz simultaneamente as coordenadas das caixas delimitadoras e as probabilidades de classe. A arquitetura divide a imagem em uma grade de tamanho $S \times S$ e, para cada célula da grade, prevê caixas delimitadoras e suas respectivas pontuações de confiança (\textit{confidence scores}).

A evolução da arquitetura YOLO ao longo dos anos (do YOLOv1 ao YOLOv8, e recentemente o YOLO11) introduziu melhorias significativas na estrutura de rede, como conexões residuais, estruturas de pirâmide de características (\textit{Feature Pyramid Networks} - FPN), mecanismos de atenção e a exclusão da dependência de âncoras (\textit{anchor-free detection}). Essas inovações otimizaram a capacidade do modelo de capturar objetos em múltiplas escalas e sob condições severas de oclusão e variação de iluminação \cite{terven2023comprehensive, ultralytics2026}.

\section{Processamento de Vídeo, Stride e Otimização com OpenVINO}
La inferência de modelos de Deep Learning em fluxos contínuos de vídeo introduz desafios relacionados à latência e ao consumo de hardware, especialmente em sistemas com hardware limitado ou CPUs sem aceleração gráfica dedicada. O processamento frame a frame de um vídeo pode gerar gargalos e provocar o esgotamento da memória RAM devido à retenção temporária de tensores em memória.

Para contornar essas limitações, utilizam-se técnicas de otimização no pipeline de processamento:
\begin{enumerate}
    \item \textbf{Amostragem Estruturada (\textit{Stride}):} A técnica de \textit{stride} na análise de vídeo consiste no pulo sistemático de frames (por exemplo, processar 1 a cada $N$ frames). Em um vídeo gravado a 30 frames por segundo (FPS), a mudança contextual de uma cena de inspeção ocorre de maneira mais lenta do que a latência de um frame isolado (33 ms). Aplicar um stride de valor 2 ou 3 reduz a carga computacional pela metade ou por um terço, mitigando problemas de timeout de requisições em APIs de borda sem prejuízo à integridade da detecção visual.
    \item \textbf{Processamento em Fluxo (\textit{Stream Processing}):} O processamento de vídeo por geradores iterativos (passando a propriedade \texttt{stream=True} na inferência do YOLO) garante que os frames sejam descartados da memória RAM imediatamente após sua análise, contendo vazamentos de memória comuns em análises em lote (\textit{batch analysis}).
    \item \textbf{OpenVINO Toolkit:} Desenvolvido pela Intel, o \textit{OpenVINO (Open Visual Inference and Neural Network Optimization)} é uma ferramenta de otimização de modelos que permite a aceleração de algoritmos de Deep Learning em plataformas de hardware Intel (CPUs comuns, gráficos integrados e GPUs dedicadas como a Intel Arc B570). O OpenVINO converte modelos originais (como PyTorch \texttt{.pt}) para uma Representação Intermediária (IR) otimizada de baixo nível, aplicando técnicas de fusão de operadores convolucionais e quantização de pesos (por exemplo, FP32 para FP16 ou INT8), o que eleva substancialmente o número de inferências por segundo (FPS) com perdas insignificantes na acurácia do modelo.
\end{enumerate}

\section{Arquitetura de APIs RESTful e Framework FastAPI}
A modularização de sistemas modernos baseia-se na separação clara entre as interfaces de usuário (aplicativos móveis, sistemas web) e as regras de negócio/processamento pesado (backend). Essa separação é tipicamente mediada por APIs (\textit{Application Programming Interfaces}) baseadas no estilo arquitetural REST (\textit{Representational State Transfer}), que preconiza a utilização de protocolo de comunicação HTTP de forma apátrida (\textit{stateless}), a identificação de recursos por meio de URLs e a manipulação de representações de dados padronizadas, normalmente em formato JSON \cite{fielding2000architectural}.

O desenvolvimento de APIs de alto desempenho em Python é beneficiado pelo framework FastAPI. O FastAPI baseia-se em conceitos modernos do Python, como suporte nativo à programação assíncrona (\texttt{async/await}), tipagem estática opcional e serialização robusta por meio da biblioteca Pydantic. Adicionalmente, o FastAPI provê autogeração de documentação interativa utilizando os padrões OpenAPI e Swagger, agilizando de forma expressiva os processos de integração entre equipes de desenvolvimento backend e mobile.

\section{Mecanismos de Segurança em APIs Web}
A segurança de dados expostos por endpoints HTTP de detecção é crítica para garantir que apenas agentes autenticados façam chamadas, evitando a sobrecarga de hardware da API por acessos não autorizados. A autenticação baseada em tokens JWT (\textit{JSON Web Tokens}) é o padrão da indústria para arquiteturas descentralizadas. O JWT armazena reivindicações (\textit{claims}) criptograficamente assinadas pelo servidor por meio de algoritmos de chave simétrica (como o HS256) ou chaves assimétricas, dispensando consultas contínuas a bancos de dados de sessão e otimizando a escalabilidade \cite{rfc7519}.

Para conferir robustez a sistemas expostos na internet, utilizam-se técnicas complementares de segurança cibernética:
\begin{itemize}
    \item \textbf{Rate Limiting:} A limitação de requisições por IP ou usuário em uma janela de tempo específica (por meio de algoritmos como \textit{sliding window}) é vital para mitigar ataques de negação de serviço (DoS) e força bruta.
    \item \textbf{404 scanning protection:} Bloqueio automático de IPs que buscam por rotas não existentes de forma sistemática (ataques de reconhecimento).
    \item \textbf{Firebase App Check e Play Integrity:} Plataformas como o Google Firebase fornecem serviços para validar se o tráfego que chega aos servidores provém exclusivamente de instâncias legítimas do aplicativo oficial e de dispositivos Android genuínos, bloqueando solicitações disparadas por bots ou emuladores modificados.
\end{itemize}

\section{Modelos de Linguagem Local (LLM) na Usabilidade de Sistemas}
Os resultados gerados por modelos de visão computacional são puramente quantitativos e numéricos (coordenadas, classes e porcentagens de confiança). Em sistemas que visam o usuário final, a apresentação dessas informações requer processamento para elevar a usabilidade e a clareza conceitual. 

A integração de Modelos de Linguagem (\textit{Large Language Models} - LLMs) locais via infraestruturas como o Ollama permite converter matrizes de detecção em relatórios descritivos textuais ricos e humanizados na língua nativa do usuário (por exemplo, português brasileiro), com total privacidade e sem latências decorrentes de conexões externas à internet ou custos de APIs proprietárias. O uso de técnicas de Engenharia de Prompts e filtros rígidos garante que o LLM comporte-se estritamente como um processador determinístico de relatórios, eliminando a ocorrência de alucinações de dados e preservando a precisão do laudo técnico emitido pelo sistema de visão.

\section{Conformidade em Segurança na Engenharia Civil}
A segurança do trabalho na indústria da construção civil é uma prioridade regulatória e humanitária. Devido à natureza mutável e fisicamente complexa dos canteiros de obras, o setor apresenta, historicamente, elevadas taxas de acidentalidade ocupacional. No Brasil, o Ministério do Trabalho e Emprego rege o setor através de Normas Regulamentadoras (NRs), dentre as quais se destacam a NR 6 \cite{brasil_nr6} (Equipamento de Proteção Individual - EPI) e a NR 18 \cite{brasil_nr18} (Segurança e Saúde no Trabalho na Indústria da Construção).

A NR 6 determina a obrigatoriedade do fornecimento gratuito e da fiscalização contínua quanto ao uso adequado de EPIs por parte do empregador. Elementos como capacetes de segurança (proteção contra impactos na cabeça), protetores auriculares (atenuação de ruídos), calçados de segurança com biqueiras de aço (proteção contra queda de materiais) e óculos de proteção (proteção contra projeção de partículas) são essenciais e legalmente exigidos para o ingresso em qualquer área do canteiro.

Já a NR 18 estipula diretrizes de conformidade para o meio físico do canteiro de obras. Isso inclui a disposição correta de equipamentos de combate a incêndio e primeiros socorros. A inspeção de conformidade visual atua como o principal instrumento de verificação destas diretrizes. A automação deste processo por meio de visão computacional fornece uma trilha de auditoria digital auditável, ajudando a constatar se a disposição física das obras está em plena concordância com as exigências técnicas nacionais.

\section{O Impacto da Segurança no Bem-estar do Trabalhador}
A segurança no trabalho estende-se para além da mera prevenção de custos corporativos ou conformidade legal, possuindo um impacto psicossocial direto e profundo sobre o bem-estar e a qualidade de vida do trabalhador da construção civil. Ambientes de trabalho desprovidos de monitoramento de riscos e com histórico de negligência quanto a normas de segurança geram altos níveis de estresse ocupacional e ansiedade crônica nos colaboradores, uma vez que a percepção de perigo iminente compromete a estabilidade emocional necessária para a execução de tarefas complexas.

A implementação de políticas de segurança ativas e o uso de sistemas modernos de monitoramento preventivo criam um ``clima de segurança'' (\textit{safety climate}) percebido de maneira positiva. Trabalhadores que atuam sob a certeza de que a empresa investe e fiscaliza a conformidade ambiental sentem-se valorizados e protegidos pelo empregador. Esse sentimento de amparo institucional está diretamente relacionado ao aumento da satisfação com o trabalho, melhoria na produtividade laboral e redução nas taxas de absenteísmo \cite{rahman2021context}. A automação e a constância nas inspeções visuais garantem que a proteção coletiva seja tratada de forma imparcial e proativa, elevando a cultura preventiva e a dignidade ocupacional do trabalhador no canteiro de obras.
